% GENERATED by backtests/spx_candle_imbalance/make_tex.py from backtests/spx_candle_imbalance/results.json
% (the one registered run of preregistration/spx_candle_imbalance.md). Do not edit numbers by hand; re-run the script.

\subsection{A test with the published imbalances}
\label{sec:imbtest}

If traders react to the published imbalance during the 15:50 minute, the candle should be right more often on days
when it points the same way as that imbalance. A pre-registered study (plan locked 2026-11-10 19:37 UTC, one registered run)
tested this prediction on the rule's 644 signal days of 2021--2026. For each year it took the 50 largest stocks
in SPY at the previous year-end, from the fund's N-PORT reports to the SEC (between 51\% and 63\% of its net assets), and bought
from Databento the imbalance messages that each stock's listing exchange sent from 15:50. From each stock's first
closing-imbalance record received in the 15:50 minute it summed the signed dollar imbalance,
$D = \sum_i s_i Q_i p_i$, with $s_i = +1$ for a buy imbalance, $-1$ for a sell imbalance and 0 for none, $Q_i$ the
imbalance in shares and $p_i$ the reference price. A day agrees if the sign of $D$ is the candle's direction. The
primary test compares the rule's hit rate on agreeing and disagreeing days, $d = h_a - h_d$, and passes if $d > 0$
and $t \ge 2.0$. By the plan's power calculation it could detect a true $d$ of about 11 to 13 points, a large concentration of
the edge on agreeing days, not a modest one. All 50 stocks had a usable record on every day, so every signal day was
classified.

\begin{table}[htbp]
\centering
\caption{Signal days that agree and disagree with the 15:50 closing imbalance, by year.}
\label{tab:imbyear}
\tabnotes{Signal days of the rule (a 15:50 body of 60\% or more), 2021--2026. A day agrees when the sign of the
summed dollar imbalance of the 50 largest stocks, from each stock's first closing-imbalance record received in the
15:50 minute, equals the candle's direction. Every signal day was classified, so no day is unclassified. Hit rate:
share of days on which the CFD moved from the 15:51 open to the 15:59 close in the candle's direction, with its
Wilson 95\% interval in percent; bp: the mean of that move in the candle's direction, in basis points. Over all
days the difference in hit rates is $-5.9$ points ($t = -1.35$), with a 20-day block-bootstrap 95\%
interval from $-14.8$ to $+4.0$ points.}
\begin{tabular}{@{}lrrlrrlr@{}}
\toprule
 & & \multicolumn{3}{c}{Agreeing days} & \multicolumn{3}{c}{Disagreeing days} \\
\cmidrule(lr){3-5}\cmidrule(l){6-8}
Year & Days & $n$ & Hit rate (95\% CI) & bp & $n$ & Hit rate (95\% CI) & bp \\
\midrule
2021 & 111 & 92 & 62.0\% (52--71) & $4.67$ & 19 & 78.9\% (57--91) & $5.12$ \\
2022 & 105 & 76 & 57.9\% (47--68) & $4.00$ & 29 & 69.0\% (51--83) & $3.88$ \\
2023 & 113 & 73 & 53.4\% (42--64) & $0.33$ & 40 & 55.0\% (40--69) & $1.18$ \\
2024 & 106 & 69 & 59.4\% (48--70) & $2.59$ & 37 & 43.2\% (29--59) & $-1.90$ \\
2025 & 117 & 91 & 53.8\% (44--64) & $1.39$ & 26 & 69.2\% (50--83) & $3.97$ \\
2026 & 92 & 68 & 47.1\% (36--59) & $0.96$ & 24 & 70.8\% (51--85) & $4.29$ \\
\midrule
All & 644 & 469 & 55.9\% (51--60) & $2.41$ & 175 & 61.7\% (54--69) & $2.24$ \\
\bottomrule
\end{tabular}
\end{table}

The candle and the imbalance usually pointed the same way, on 469 of the 644 days
(72.8\%). But agreement did not make the candle more accurate (Table~\ref{tab:imbyear}): it was right on
55.9\% of agreeing days and 61.7\% of disagreeing days, $d = -5.9$
points ($t = -1.35$), with a block-bootstrap interval from $-14.8$ to $+4.0$ points,
and disagreeing days had the higher hit rate in five of the six years, all but 2024. The imbalance direction alone was right on 51.1\% of days. The robustness checks fixed in
the plan agree: the difference in hit rates stayed negative with the latest record of the minute, with NYSE or
Nasdaq stocks alone and within year and direction, and the differences in the average move and in the option P\&L
were small against their noise ($t = 0.13$ and $0.79$; Internet Appendix). The prediction was
therefore not confirmed. The test has limited power, and a failure does not refute the auction explanation: the
price may adjust within the 15:50 candle itself, imbalances change after their first publication, and 50 stocks are
only part of the index. The high agreement is an association, not a cause, since the candle and the imbalance can
reflect the same closing orders. What the test does show is that the rule's edge was not concentrated on the days on
which its candle agreed with the imbalance first published at 15:50.
